\documentclass[runningheads]{llncs}
\usepackage[T1]{fontenc}
\usepackage[utf8]{inputenc}
\usepackage{graphicx}
\usepackage{booktabs}
\usepackage{amsmath}
\usepackage{amssymb}
\usepackage{cite}
\usepackage{hyperref}
\usepackage{xcolor}
\usepackage{xspace}

\graphicspath{{figures/}}

\begin{document}

\title{Compilation-Free Smart Contract Vulnerability Detection\\via Static Feature Engineering and Ensemble Learning}

\author{Author information omitted for double-blind review}

\maketitle

% ===========================================================================
\begin{abstract}
Detecting vulnerabilities in smart contracts is critical for securing decentralized applications, yet existing analysis tools require successful compilation with specific Solidity compiler versions---a barrier that limits their applicability to large-scale, heterogeneous contract corpora.
We present a compilation-free approach that extracts 292 static features directly from Solidity source text using regex-based parsing and lightweight AST analysis.
Our feature set spans six categories: code complexity metrics, Halstead software science measures, call-graph approximations, data-flow indicators, vulnerability-pattern templates, and security-practice signals.
We evaluate the approach on 26,423 contracts from three public datasets (DAppSCAN, FORGE, and AccessControlVulnerabilities) using GroupKFold cross-validation to prevent data leakage across projects.
On the DAppSCAN dataset alone, a GradientBoosting classifier achieves ROC-AUC\,=\,0.894, PR-AUC\,=\,0.363, and F1\,=\,0.400 with 4.4\% class prevalence.
By merging DAppSCAN and FORGE for training, performance improves substantially to ROC-AUC\,=\,0.920, PR-AUC\,=\,0.497, and F1\,=\,0.592---a 37\% relative gain in PR-AUC and 48\% in F1.
A two-stage feature selection procedure (mutual information filtering followed by Random Forest importance ranking) reduces the feature space from 292 to 80 features with no loss in discriminative power, and cross-dataset transfer experiments confirm generalization: a DAppSCAN-trained model achieves 99.5--100\% recall on FORGE vulnerabilities and 100\% recall on all 94 AccessControlVulnerabilities contracts.
The results demonstrate that static source-code features alone, without compilation or bytecode analysis, can effectively prioritize smart contracts for security auditing.
\end{abstract}

\keywords{Smart contract security \and Vulnerability detection \and Static analysis \and Feature engineering \and Machine learning \and Solidity}

% ===========================================================================
\section{Introduction}
\label{sec:intro}

Smart contracts deployed on Ethereum and other blockchain platforms manage billions of dollars in digital assets, yet remain susceptible to exploitable vulnerabilities.
High-profile incidents such as the DAO hack (\$60M, 2016)~\cite{mehar2019understanding}, the Parity wallet freeze (\$280M, 2017)~\cite{destefanis2018smart}, and the Ronin bridge exploit (\$625M, 2022)~\cite{zhou2023sok} underscore the critical need for automated vulnerability detection.
Once deployed, smart contracts are immutable---bugs cannot be patched without costly proxy-upgrade mechanisms or redeployment---making pre-deployment detection essential.

Existing automated analysis tools fall into two broad categories.
\emph{Static analysis} tools such as Slither~\cite{feist2019slither} and SmartCheck~\cite{tikhomirov2018smartcheck} detect vulnerability patterns through predefined rules over the abstract syntax tree (AST) or intermediate representations.
\emph{Symbolic execution} tools such as Mythril~\cite{mueller2018smashing} and Oyente~\cite{luu2016making} explore execution paths to discover violations of safety properties.
While these tools achieve high precision for the specific vulnerability types they target, empirical evaluations reveal significant limitations: a SmartBugs benchmark of nine tools on 47,587 contracts found that only 42\% of annotated vulnerabilities were detected, with 97\% of contracts flagged as vulnerable---indicating severe false-positive rates~\cite{durieux2020empirical}.
Moreover, all of these tools require successful compilation with a compatible Solidity compiler version, which is a non-trivial barrier for large-scale analysis: contracts span Solidity versions 0.4.x through 0.8.x, often depend on external libraries, and may use non-standard import paths.

Machine-learning approaches have been proposed as an alternative, using graph neural networks (GNNs) on contract control-flow graphs~\cite{zhuang2020smart}, code embeddings~\cite{gao2020smartembed}, or sequence models on opcode~\cite{lutz2021escort}.
However, these methods typically require compilation to obtain bytecode, control-flow graphs, or ASTs from the Solidity compiler, inheriting the same compilation barrier.

In this paper, we take a different approach: we extract a comprehensive set of 292 static features \emph{directly from Solidity source text} using regex-based parsing and a lightweight, non-compiling AST parser.
No Solidity compiler, bytecode, or deployment to a blockchain is required.
Our features span six categories designed to capture complementary aspects of contract risk:

\begin{enumerate}
    \item \textbf{Code complexity metrics:} cyclomatic complexity (per-function and aggregated), lines of code, function count and size distribution, nesting depth, and structural concentration measures (Gini coefficient, coefficient of variation).
    \item \textbf{Halstead software science:} vocabulary size, volume, difficulty, and effort---established software engineering metrics that quantify cognitive complexity.
    \item \textbf{Call-graph approximations:} edge count, density, in/out degree, leaf-function ratio, and maximum call depth---extracted via regex matching of function declarations and invocations.
    \item \textbf{Data-flow indicators:} state-variable declaration/write counts, write fanout entropy, shared-write ratios, and read-after-write pair counts.
    \item \textbf{Vulnerability-pattern templates:} counts of reentrancy-prone sites, unchecked low-level calls, \texttt{delegatecall} usage, \texttt{tx.origin} authentication, and missing zero-address checks.
    \item \textbf{Security-practice signals:} Checks-Effects-Interactions (CEI) compliance ratio, access-control coverage, event emission on state changes, \texttt{require}/\texttt{assert} density, and comment-to-code ratio.
\end{enumerate}

We evaluate this feature set on 26,423 contracts drawn from three publicly available datasets: DAppSCAN~\cite{wu2023dappscan} (21,457 contracts with clean/vulnerable labels from 1,139 audited DApp projects), FORGE (4,872 contracts with vulnerable/suspicious/unknown labels), and AccessControlVulnerabilities (94 contracts with confirmed access-control vulnerabilities).
Using rigorous GroupKFold cross-validation---where folds are split by project to prevent data leakage---we train and evaluate seven model families, from logistic regression to gradient-boosted ensembles.

\subsection{Contributions}

The main contributions of this work are:

\begin{itemize}
    \item \textbf{Compilation-free feature engineering.} We design and implement 292 static features extractable from raw Solidity source text without compilation, spanning six complementary categories. A two-stage feature selection procedure identifies 80 features that retain full discriminative power.
    
    \item \textbf{Large-scale, multi-dataset evaluation.} We evaluate on 26,423 contracts across three datasets using project-level GroupKFold cross-validation---a stricter and more realistic protocol than the random train/test splits commonly used in the literature.
    
    \item \textbf{Cross-dataset generalization analysis.} We present four cross-dataset experiments (single-dataset CV, cross-prediction, reverse transfer, and merged training) demonstrating that the feature set generalizes across dataset boundaries, with merged training achieving ROC-AUC\,=\,0.920 and F1\,=\,0.592.
    
    \item \textbf{Interpretable risk assessment.} The top features---composite complexity, risky code bigrams, state-variable write counts, and Halstead effort---are directly interpretable by auditors and align with known vulnerability drivers in Solidity.
\end{itemize}

The remainder of the paper is organized as follows.
Section~\ref{sec:related} reviews related work on smart contract vulnerability detection.
Section~\ref{sec:method} describes the feature engineering methodology and machine learning pipeline.
Section~\ref{sec:data} presents the datasets used for evaluation.
Section~\ref{sec:results} reports experimental results, including single-dataset, cross-dataset, and ablation studies.
Section~\ref{sec:discussion} discusses findings, limitations, and practical implications.
Section~\ref{sec:conclusion} concludes the paper.

% ===========================================================================
\section{Related Work}
\label{sec:related}

\subsection{Rule-Based and Symbolic Analysis Tools}

The most widely adopted approaches for smart contract vulnerability detection rely on predefined rules or formal methods applied to compiled contract representations.

\textbf{Static analysis.}
Slither~\cite{feist2019slither} performs static analysis on the Solidity compiler's intermediate representation (SlithIR), offering over 80 built-in vulnerability detectors with fast execution (seconds per contract).
SmartCheck~\cite{tikhomirov2018smartcheck} translates Solidity source into an XML parse tree and applies XPath patterns to detect anti-patterns.
Securify~\cite{tsankov2018securify} extracts semantic information from EVM bytecode to check compliance and violation patterns against security properties.

\textbf{Symbolic execution.}
Mythril~\cite{mueller2018smashing} uses symbolic execution and concolic analysis to explore execution paths and detect vulnerabilities such as integer overflows and reentrancy.
Oyente~\cite{luu2016making} was among the first tools to apply symbolic execution to Ethereum contracts, targeting four vulnerability types.
Manticore~\cite{mossberg2019manticore} combines symbolic execution with concrete inputs for vulnerability discovery.

\textbf{Formal verification.}
VeriSolid~\cite{mavridou2019verisolid} applies model checking to finite-state machine representations of contracts.
VerX~\cite{permenev2020verx} verifies functional properties of Ethereum contracts via delayed abstraction.

\textbf{Empirical benchmarking.}
Large-scale evaluations have revealed significant limitations.
Durieux et al.~\cite{durieux2020empirical} evaluated nine tools on 47,587 Ethereum contracts and found only 42\% vulnerability detection rate with a 97\% false-positive rate.
The SolidiFI framework~\cite{ghaleb2020effective} injected 9,369 bugs into 50 contracts and found that all six evaluated tools (including Slither and Mythril) exhibited both false negatives and false positives.
The DAppSCAN benchmark~\cite{wu2023dappscan} tested five state-of-the-art tools on real-world audit data and reported ``subpar performance in terms of both effectiveness and success detection rate.''

A fundamental limitation shared by all these tools is the requirement for successful compilation.
Contracts must be compiled with a compatible Solidity compiler version, and external dependencies must be resolved---a non-trivial constraint when analyzing heterogeneous corpora spanning multiple years and Solidity versions.

\subsection{Machine Learning Approaches}

Several machine learning methods have been proposed for vulnerability detection.

\textbf{Graph-based methods.}
Zhuang et al.~\cite{zhuang2020smart} proposed DR-GCN and TMP, which apply graph neural networks to contract control-flow and data-flow graphs to detect reentrancy and timestamp dependence.
Liu et al.~\cite{liu2021combining} combined expert patterns with GNNs for improved detection.
These methods require compilation to construct the graph representations.

\textbf{Sequence-based methods.}
Escort~\cite{lutz2021escort} uses a multi-label deep neural network on EVM opcode sequences.
ContractWard~\cite{wang2020contractward} extracts bigram features from opcode and applies machine learning classifiers, achieving strong results but requiring bytecode.
Eth2Vec~\cite{ashizawa2021eth2vec} learns vector representations of EVM bytecode for vulnerability classification.

\textbf{Source-level methods.}
SmartEmbed~\cite{gao2020smartembed} uses code embeddings from the AST for clone and bug detection.
Peculiar~\cite{wu2021peculiar} combines critical program paths with attention-based neural networks.
These methods typically still require a compiler-generated AST.

\textbf{Gap addressed by our work.}
To the best of our knowledge, no prior work has proposed a comprehensive, compilation-free feature set for contract-level vulnerability classification.
Our approach differs from existing ML methods in three key ways: (1)~it requires no compilation, bytecode, or AST from the Solidity compiler; (2)~it operates at the contract level rather than targeting specific vulnerability types; and (3)~it produces both a continuous risk score for prioritization and a binary classification for triage.

% ===========================================================================
\section{Methodology}
\label{sec:method}

Our approach consists of three stages: (1)~feature extraction from raw Solidity source text, (2)~feature selection to identify the most discriminative subset, and (3)~classification using supervised machine learning models with rigorous cross-validation.

\subsection{Feature Extraction}
\label{sec:features}

We extract 292 numerical features from each contract's Solidity source code using regex-based parsing and a lightweight non-compiling AST parser (\texttt{solidity\_parser}).
No Solidity compiler invocation is required.
The features are organized into six categories, summarized in Table~\ref{tab:feature_categories}.

\begin{table}[t]
\centering
\caption{Feature categories and counts. All features are extracted from raw Solidity source text without compilation.}
\label{tab:feature_categories}
\begin{tabular}{lrl}
\hline
\textbf{Category} & \textbf{Count} & \textbf{Examples} \\
\hline
Code complexity       &  68 & Cyclomatic complexity (avg, std, max), LOC, \\
                      &     & function count/size distribution, nesting depth \\
Halstead metrics      &   4 & Volume, difficulty, effort, volume per LOC \\
Token \& bigram       &  18 & Token entropy/diversity, risky bigrams, \\
                      &     & DeFi bigrams, bigram entropy \\
State \& data flow    &  24 & State vars written, fanout entropy, \\
                      &     & shared-write ratio, read-after-write pairs \\
Call graph            &   7 & Edges, density, max degree, leaf ratio, depth \\
Vulnerability patterns&  18 & Reentrancy sites, unchecked calls, \\
\& security practices &     & CEI compliance, access-control coverage \\
Function-level        &  46 & Per-function LOC/arithmetic/assignment \\
aggregations          &     & statistics (mean, std, Gini, percentiles) \\
Risk-score components &   9 & Composite scores for complexity, safety, \\
                      &     & interaction, privilege, code smell \\
Other structural      &  98 & Imports, inheritance, type casts, \\
                      &     & comment ratio, compiler features \\
\hline
\textbf{Total}        & \textbf{292} & \\
\hline
\end{tabular}
\end{table}

\textbf{Code complexity metrics.}
For each function in the contract, we compute cyclomatic complexity by counting control-flow branches (\texttt{if}, \texttt{for}, \texttt{while}, \texttt{require}, logical operators) plus one for the base path.
We aggregate per-function complexities into max, mean, standard deviation, sum, per-LOC density, and variance ratio.
Function-size distributions are characterized via coefficient of variation, Gini coefficient, inter-quartile range, and skewness.

\textbf{Halstead software science metrics.}
We identify operators (symbols such as \texttt{+}, \texttt{-}, \texttt{*}, \texttt{=}) and operands (identifiers and numeric literals) via regex and compute the standard Halstead metrics: vocabulary $n = n_1 + n_2$, length $N = N_1 + N_2$, volume $V = N \log_2 n$, difficulty $D = \frac{n_1}{2} \cdot \frac{N_2}{n_2}$, and effort $E = V \cdot D$.

\textbf{Call-graph approximations.}
We extract function declarations and function-call sites via regex, build an approximate call graph, and compute edge count, density, maximum in-degree and out-degree, leaf-function ratio, and maximum call depth via iterative traversal.

\textbf{State and data-flow indicators.}
State variables are identified from declarations matching Solidity type patterns.
We count how many are actually written to (assigned) anywhere in the contract, compute the state-variable write fanout (how many functions access each variable) and its Shannon entropy, and detect shared-write variables (written by multiple functions) and read-after-write pairs.

\textbf{Vulnerability-pattern templates.}
We count occurrences of known vulnerability patterns using multi-line regex: reentrancy-prone sites (external calls followed by state writes), unchecked low-level calls (\texttt{.call\{\}} without return-value checks), \texttt{delegatecall} usage, \texttt{tx.origin} in conditional expressions (phishing vulnerability), and missing zero-address checks before \texttt{transfer}.

\textbf{Security-practice signals.}
We measure positive security practices: the ratio of functions following the Checks-Effects-Interactions (CEI) pattern, the fraction of state-modifying functions with access-control modifiers, event emission on state changes, \texttt{require}/\texttt{assert} density per function, and the comment-to-code ratio.

\subsection{Feature Selection}
\label{sec:feature_selection}

To identify the most informative features and reduce dimensionality, we apply a two-stage selection procedure:

\begin{enumerate}
    \item \textbf{Mutual information filter.} We compute the mutual information (MI) between each feature and the binary vulnerability label.  Features with zero MI (no statistical association with the label) are removed.  This typically reduces the feature space from 292 to approximately 240 features.
    
    \item \textbf{Random Forest importance ranking.} We fit a Random Forest classifier on the MI-filtered features and rank by mean decrease in impurity (Gini importance).  The top $k = 80$ features are retained.
\end{enumerate}

This procedure is applied independently within each cross-validation fold to prevent information leakage from the test set into feature selection.

\subsection{Classification Models}
\label{sec:models}

We evaluate seven classifier families spanning linear, ensemble, and boosting approaches:

\begin{itemize}
    \item \textbf{Logistic Regression} (L2 and L1 regularization, balanced class weights)
    \item \textbf{Random Forest} (500 trees, balanced subsampling)
    \item \textbf{Extra Trees} (500 trees, balanced subsampling)
    \item \textbf{GradientBoosting} (300 estimators, learning rate 0.05, subsampling 0.8)
    \item \textbf{HistGradientBoosting} (500 iterations, early stopping, balanced class weights)
    \item \textbf{XGBoost} (500 estimators, automatic \texttt{scale\_pos\_weight})
\end{itemize}

All models use class-weight balancing or equivalent mechanisms to handle the severe class imbalance (4.4\% positive rate in DAppSCAN).
Features are standardized (zero mean, unit variance) before training.

\subsection{Evaluation Protocol}
\label{sec:eval_protocol}

We use 5-fold GroupKFold cross-validation, where groups correspond to project identifiers.
This ensures that all contracts from the same audited project appear in the same fold, preventing the model from memorizing project-specific patterns and leaking information across train/test boundaries.

For each fold, the model is trained on the training set and evaluated on the held-out test set.
We compute:
\begin{itemize}
    \item \textbf{ROC-AUC}: area under the receiver operating characteristic curve (threshold-free ranking quality).
    \item \textbf{PR-AUC}: area under the precision-recall curve (more informative than ROC-AUC under class imbalance).
    \item \textbf{F1 score}: computed at the threshold that maximizes F1 on the test fold.
    \item \textbf{Precision and Recall}: at the optimal F1 threshold.
\end{itemize}

We report mean and standard deviation across the five folds for ROC-AUC and PR-AUC, and aggregate confusion matrix entries (TP, FP, FN, TN) across all folds.

As a baseline, we evaluate the rule-based composite risk score (\texttt{risk\_z}, a weighted sum of nine risk components) using the same GroupKFold protocol.

% ===========================================================================
\section{Datasets}
\label{sec:data}

\begin{table}[t]
\centering
\caption{Overview of evaluation datasets.}
\label{tab:datasets}
\begin{tabular}{lrrrrr}
\hline
\textbf{Dataset} & \textbf{Contracts} & \textbf{Clean} & \textbf{Vulnerable} & \textbf{Other} & \textbf{Pos.\ Rate} \\
\hline
DAppSCAN & 21{,}457 & 20{,}509 & 948 & 0 & 4.4\% \\
FORGE & 4{,}872 & 0 & 1{,}244 & 3{,}628$^\dagger$ & 52.6\%$^\ddagger$ \\
AccessControlVuln & 94 & 0 & 94 & 0 & 100\% \\
\hline
\textbf{Total} & \textbf{26{,}423} & 20{,}509 & 2{,}286 & 3{,}628 & --- \\
\hline
\end{tabular}
\\[4pt]
{\footnotesize $^\dagger$2{,}505 suspicious + 1{,}123 unknown. $^\ddagger$Vulnerable vs.\ unknown only.}
\end{table}


We evaluate on three publicly available datasets comprising 26,423 smart contracts in total (Table~\ref{tab:datasets}).

\textbf{DAppSCAN}~\cite{wu2023dappscan} is the largest and best-labeled dataset, containing 21,457 Solidity source files from 1,139 real-world DApp projects that underwent professional security audits.
Each contract is labeled as either \emph{clean} (20,509) or \emph{vulnerable} (948), yielding a positive rate of 4.4\%.
The severe class imbalance reflects real-world prevalence---most audited contracts pass review---and makes PR-AUC the most informative evaluation metric.

\textbf{FORGE} contains 4,872 contracts labeled as \emph{vulnerable} (1,244), \emph{suspicious} (2,505), or \emph{unknown} (1,123).
The absence of a \emph{clean} label limits its standalone use for binary classification.
For cross-dataset experiments, we treat \emph{vulnerable} as positive and \emph{unknown} as negative (dropping the ambiguous \emph{suspicious} class), yielding 2,367 contracts with a 52.6\% positive rate.

\textbf{AccessControlVulnerabilities} is a curated collection of 94 contracts with confirmed access-control vulnerabilities.
With no negative examples, it serves as a sanity check: a well-trained model should flag all 94 contracts as vulnerable.

% ===========================================================================
\section{Results}
\label{sec:results}

We present results in four parts: single-dataset evaluation on DAppSCAN (Section~\ref{sec:single}), cross-dataset experiments (Section~\ref{sec:cross}), feature analysis (Section~\ref{sec:feat_analysis}), and ablation studies (Section~\ref{sec:ablation}).

\subsection{Single-Dataset Evaluation (DAppSCAN)}
\label{sec:single}

\begin{table}[t]
\centering
\caption{Classification performance on DAppSCAN (single-dataset) and merged (DAppSCAN + FORGE) training. Best values in \textbf{bold}.}
\label{tab:model_comparison}
\setlength{\tabcolsep}{4pt}
\begin{tabular}{ll ccccc}
\hline
\textbf{Training} & \textbf{Model} & \textbf{ROC-AUC} & \textbf{PR-AUC} & \textbf{F1} & \textbf{Prec.} & \textbf{Recall} \\
\hline
DAppSCAN & Random Forest & 0.901$\pm$0.006 & 0.307 & 0.393 & 0.326 & 0.500 \\
DAppSCAN & XGBoost & 0.897$\pm$0.010 & 0.348 & 0.393 & 0.333 & 0.492 \\
DAppSCAN & GradientBoosting & 0.894$\pm$0.009 & 0.363 & 0.400 & 0.341 & 0.505 \\
DAppSCAN & Extra Trees & 0.894$\pm$0.007 & 0.302 & 0.378 & 0.296 & 0.525 \\
DAppSCAN & HistGradientBoosting & 0.887$\pm$0.008 & 0.302 & 0.382 & 0.324 & 0.485 \\
DAppSCAN & Logistic Reg. & 0.856$\pm$0.006 & 0.256 & 0.338 & 0.303 & 0.414 \\
DAppSCAN & Baseline (risk score) & 0.815$\pm$0.015 & 0.168 & 0.255 & 0.170 & 0.518 \\
\hline
Merged & Random Forest & \textbf{0.920}$\pm$0.006 & 0.478 & 0.579 & 0.482 & \textbf{0.729} \\
Merged & GradientBoosting & \textbf{0.920}$\pm$0.006 & \textbf{0.497} & \textbf{0.592} & 0.500 & 0.726 \\
Merged & XGBoost & 0.918$\pm$0.006 & 0.491 & 0.582 & 0.503 & 0.696 \\
Merged & HistGradientBoosting & 0.916$\pm$0.005 & 0.476 & 0.575 & \textbf{0.510} & 0.662 \\
\hline
\end{tabular}
\end{table}

Table~\ref{tab:model_comparison} reports classification performance on DAppSCAN using 5-fold GroupKFold cross-validation.
All ML models substantially outperform the rule-based baseline (ROC-AUC\,=\,0.815), with ensemble and boosting models achieving ROC-AUC above 0.88.

Among individual models, Random Forest achieves the highest ROC-AUC (0.901), while GradientBoosting achieves the best PR-AUC (0.363) and F1 (0.400).
The strong performance of boosting methods on PR-AUC suggests they are better calibrated for the severe class imbalance.

The merged training configuration (DAppSCAN + FORGE) yields dramatic improvements across all metrics and all models.
GradientBoosting on the merged data achieves ROC-AUC\,=\,0.920, PR-AUC\,=\,0.497, and F1\,=\,0.592, representing relative improvements of 2.9\%, 37.1\%, and 48.1\% over the DAppSCAN-only configuration, respectively.

\begin{figure}[t]
    \centering
    \includegraphics[width=\textwidth]{fig_model_comparison.pdf}
    \caption{Model performance comparison: DAppSCAN-only vs.\ merged (DAppSCAN + FORGE) training across three metrics. Error bars show standard deviation across five folds.}
    \label{fig:model_comparison}
\end{figure}

Figure~\ref{fig:model_comparison} visualizes the improvement from merged training.
The gain is most pronounced for PR-AUC and F1, confirming that the additional FORGE data helps the model better distinguish vulnerable contracts in the precision-recall space where class imbalance matters most.

\subsection{Cross-Dataset Experiments}
\label{sec:cross}

\begin{table}[t]
\centering
\caption{Cross-dataset transfer results (GradientBoosting). Models trained on one dataset and evaluated on another.}
\label{tab:cross_dataset}
\resizebox{\textwidth}{!}{%
\begin{tabular}{ll ccccc r}
\hline
\textbf{Train On} & \textbf{Test On} & \textbf{ROC-AUC} & \textbf{PR-AUC} & \textbf{F1} & \textbf{Prec.} & \textbf{Recall} & \textbf{N} \\
\hline
DAppSCAN & DAppSCAN (CV) & 0.895 & 0.369 & 0.416 & 0.365 & 0.487 & 21,457 \\
DAppSCAN & FORGE & 0.492 & 0.527 & 0.689 & 0.525 & 0.999 & 2,367 \\
DAppSCAN & AccessControlVuln & ---$^{\dagger}$ & ---$^{\dagger}$ & 1.000$^{\dagger}$ & 1.000$^{\dagger}$ & 1.000 & 94 \\
FORGE & DAppSCAN & 0.643 & 0.075 & 0.151 & 0.099 & 0.314 & 21,457 \\
FORGE & FORGE (CV) & 0.553 & 0.561 & 0.703 & 0.546 & 0.986 & 2,367 \\
DAppSCAN+FORGE & Merged (CV) & 0.921 & 0.506 & 0.597 & 0.503 & 0.734 & 23,824 \\
\hline
\end{tabular}}
\vspace{4pt}

{\footnotesize $^{\dagger}$AccessControlVuln contains only positive samples (94 vulnerable, 0 clean). ROC-AUC and PR-AUC are undefined (require both classes). F1 and Precision are trivially 1.0 since FP\,=\,0.}
\end{table}

\begin{figure}[t]
    \centering
    \includegraphics[width=\textwidth]{fig_cross_dataset_heatmap.pdf}
    \caption{Cross-dataset transfer performance (GradientBoosting). Rows indicate training data; columns indicate test data. Left: ROC-AUC. Right: F1 score.}
    \label{fig:cross_heatmap}
\end{figure}

Table~\ref{tab:cross_dataset} and Figure~\ref{fig:cross_heatmap} summarize four cross-dataset experiments.

\textbf{DAppSCAN $\rightarrow$ FORGE \& AccessControl.}
A model trained exclusively on DAppSCAN achieves near-perfect recall (99.9\%) on FORGE's vulnerable contracts and 100\% recall on all 94 AccessControlVulnerabilities contracts.
However, ROC-AUC on FORGE is only 0.492, indicating the model cannot distinguish FORGE's \emph{vulnerable} from \emph{unknown} contracts---it flags both as risky.
This is consistent with the hypothesis that FORGE's \emph{unknown} contracts may themselves be partially vulnerable.

\textbf{FORGE $\rightarrow$ DAppSCAN.}
The reverse transfer (training on FORGE, predicting on DAppSCAN) achieves only ROC-AUC\,=\,0.643 and F1\,=\,0.151.
FORGE's small size (2,367 contracts) and lack of clean examples make it insufficient for learning a generalizable decision boundary.

\textbf{Merged training.}
Combining both datasets for GroupKFold CV yields the strongest results (ROC-AUC\,=\,0.920, PR-AUC\,=\,0.497, F1\,=\,0.592), confirming that the datasets contribute complementary training signal.

\begin{figure}[t]
    \centering
    \includegraphics[width=\textwidth]{fig_confusion_matrices.pdf}
    \caption{Confusion matrices for the best model (GradientBoosting). Left: DAppSCAN-only training. Right: merged training. Merged training increases TP from 454 to 1,591 while maintaining a low false-positive rate (6.7\%).}
    \label{fig:confusion}
\end{figure}

Figure~\ref{fig:confusion} compares confusion matrices.
Merged training increases the number of true positives from 454 to 1,591 (recall from 47.7\% to 72.6\%) while the false-positive rate increases only modestly from 3.9\% to 6.7\%.

\subsection{Feature Analysis}
\label{sec:feat_analysis}

\begin{table}[t]
\centering
\caption{Top 20 features by importance (two-stage MI + RF selection, 257$\rightarrow$80).}
\label{tab:features}
\begin{tabular}{rlll}
\hline
\textbf{Rank} & \textbf{Feature} & \textbf{Description} & \textbf{Importance} \\
\hline
1 & \texttt{comp\_complexity} & Composite complexity score & 0.0409 \\
2 & \texttt{n\_risky\_bigrams} & Risky code bigram count & 0.0304 \\
3 & \texttt{n\_written\_state\_vars} & State variables written & 0.0294 \\
4 & \texttt{halstead\_effort} & Halstead effort metric & 0.0261 \\
5 & \texttt{halstead\_difficulty} & Halstead difficulty metric & 0.0257 \\
6 & \texttt{n\_assignments} & Assignment statement count & 0.0251 \\
7 & \texttt{n\_arithmetic\_ops} & Arithmetic operation count & 0.0204 \\
8 & \texttt{loc} & Lines of code & 0.0202 \\
9 & \texttt{std\_cyclomatic} & Cyclomatic complexity (std) & 0.0197 \\
10 & \texttt{bigram\_entropy} & Code bigram entropy & 0.0194 \\
11 & \texttt{risky\_bigram\_density} & Risky bigram density & 0.0171 \\
12 & \texttt{sv\_fanout\_entropy} & State-variable fanout entropy & 0.0157 \\
13 & \texttt{fn\_loc\_gini} & Function size Gini coefficient & 0.0153 \\
14 & \texttt{cyclomatic\_variance\_ratio} & Cyclomatic variance ratio & 0.0148 \\
15 & \texttt{arithmetic\_density} & Arithmetic density & 0.0146 \\
16 & \texttt{fn\_loc\_cv} & Function size coeff. of variation & 0.0134 \\
17 & \texttt{n\_state\_variables} & State variable count & 0.0132 \\
18 & \texttt{avg\_cyclomatic} & Avg. cyclomatic complexity & 0.0128 \\
19 & \texttt{std\_fn\_arith} & Arithmetic per-function (std) & 0.0126 \\
20 & \texttt{token\_entropy} & Token entropy & 0.0116 \\
\hline
\end{tabular}
\end{table}

\begin{figure}[t]
    \centering
    \includegraphics[width=0.7\textwidth]{fig_feature_importance.pdf}
    \caption{Top 20 features by importance (two-stage MI + RF selection, 292 $\rightarrow$ 80).}
    \label{fig:feat_imp}
\end{figure}

Table~\ref{tab:features} and Figure~\ref{fig:feat_imp} show the top 20 features after two-stage selection.
The most important feature is \texttt{comp\_complexity} (composite complexity score), followed by \texttt{n\_risky\_bigrams} (count of dangerous token pairs such as \texttt{msg.sender}, \texttt{msg.value}, \texttt{tx.origin}), and \texttt{n\_written\_state\_vars} (number of state variables actually assigned to).

The top features cluster into four interpretable groups:
(1)~\emph{complexity metrics} (composite complexity, Halstead effort/difficulty, cyclomatic std),
(2)~\emph{state mutation indicators} (written state vars, assignments, state-variable fanout),
(3)~\emph{dangerous code patterns} (risky bigrams, arithmetic operations and density), and
(4)~\emph{code structure} (LOC, function-size Gini, token entropy).
All of these align with known vulnerability drivers in Solidity: complex contracts with extensive state mutation and frequent use of sensitive values (\texttt{msg.value}, \texttt{block.timestamp}) present larger attack surfaces.

\subsection{Ablation Studies}
\label{sec:ablation}

\begin{figure}[t]
    \centering
    \includegraphics[width=\textwidth]{fig_feature_ablation.pdf}
    \caption{Feature set ablation on DAppSCAN. ``All'' includes 292 features (static + on-chain + bytecode); ``Static Only'' uses 257 source-code features; ``Components Only'' uses 9 risk-score components.}
    \label{fig:ablation}
\end{figure}

We evaluate three feature subsets (Figure~\ref{fig:ablation}):

\textbf{All features (292).}
The full feature set including 257 static source-code features, 24 on-chain activity features, and 11 bytecode features.

\textbf{Static only (257).}
Source-code features only, excluding on-chain and bytecode features.
ROC-AUC is virtually unchanged (0.901 vs.\ 0.901 for Random Forest), confirming that the on-chain and bytecode features---which are nearly all zero in this dataset---contribute negligibly.
This demonstrates that our approach works purely from source code.

\textbf{Components only (9).}
Only the nine high-level risk-score components (complexity, safety, interaction, privilege, code smell, etc.).
ROC-AUC drops to 0.856 but remains well above the baseline (0.815), showing that even a compact 9-feature summary retains substantial discriminative power.

\textbf{Feature selection (80 of 292).}
The two-stage selection procedure (Section~\ref{sec:feature_selection}) identifies 80 features that achieve ROC-AUC\,=\,0.903 (Random Forest) and F1\,=\,0.400---matching or slightly exceeding the full 292-feature set.
This confirms that approximately two-thirds of features are redundant and that a compact, interpretable feature set suffices.

% ===========================================================================
\section{Discussion}
\label{sec:discussion}

\subsection{Comparison with Existing Approaches}

Our compilation-free method occupies a distinct niche in the smart contract analysis landscape.
Traditional tools (Slither, Mythril) detect specific vulnerability types with high precision but require compilation and have limited recall on real-world benchmarks~\cite{durieux2020empirical,ghaleb2020effective,wu2023dappscan}.
Our method instead provides a holistic, contract-level risk assessment without compilation.

The approaches are complementary: our method can serve as a first-pass screening tool to prioritize contracts for more expensive analysis by Slither, Mythril, or manual audit.
With a recall of 72.6\% (merged training) and precision of 50.0\%, approximately half of flagged contracts are truly vulnerable, and nearly three-quarters of vulnerable contracts are caught---a useful triage ratio for large-scale screening.

\subsection{Limitations}

\textbf{Semantic vulnerabilities.}
Features based on structural patterns cannot capture purely semantic vulnerabilities such as business-logic errors or incorrect access-control policies that are syntactically indistinguishable from correct code.
This places a ceiling on recall achievable by any static-feature approach.

\textbf{Label quality.}
DAppSCAN labels are derived from professional audit reports, which may contain false negatives (vulnerabilities missed by auditors) or inconsistent severity thresholds across audit teams.
FORGE's \emph{suspicious} and \emph{unknown} categories introduce additional label noise.

\textbf{Dataset coverage.}
The evaluated contracts primarily come from DeFi projects audited between 2019--2023.
Generalization to other domains (NFT, gaming, cross-chain bridges) or newer Solidity patterns requires further validation.

\textbf{Granularity.}
Our method classifies at the contract level (vulnerable vs.\ clean) rather than identifying specific vulnerability types or locations within the code.
Combining our approach with targeted tools like Slither could provide both prioritization and precise localization.

\subsection{Practical Implications}

The compilation-free nature of our approach enables several practical applications:
(1)~\emph{Large-scale screening} of contract repositories without managing compiler dependencies;
(2)~\emph{CI/CD integration} where a risk score is computed on every commit before attempting compilation;
(3)~\emph{Audit prioritization} where the continuous risk score ranks contracts for manual review;
(4)~\emph{Historical analysis} of contracts spanning multiple Solidity versions without version-specific toolchains.

% ===========================================================================
\section{Conclusion}
\label{sec:conclusion}

We presented a compilation-free approach to smart contract vulnerability detection based on 292 static features extracted directly from Solidity source text.
Evaluated on 26,423 contracts across three datasets with rigorous GroupKFold cross-validation, our best configuration (merged DAppSCAN + FORGE training with GradientBoosting) achieves ROC-AUC\,=\,0.920, PR-AUC\,=\,0.497, and F1\,=\,0.592.
Cross-dataset transfer experiments confirm generalization, with 100\% recall on known-vulnerable AccessControlVulnerabilities contracts.
Feature analysis reveals that code complexity, state mutation patterns, and dangerous code bigrams are the strongest predictors, all directly interpretable by security auditors.

Future work includes extending the feature set with lightweight semantic analysis (e.g., taint propagation via regex), evaluating on contract families beyond DeFi, and combining the continuous risk score with targeted vulnerability detectors for a two-stage audit pipeline.

% ===========================================================================
\bibliographystyle{splncs04}
\bibliography{sc_vulnerability_refs}

\end{document}
